\section{The Autonomous Drift Problem}
\label{sec:drift}

Before specifying the Recursive Spawning Protocol, we must precisely characterize the failure mode it targets: \emph{autonomous drift} — the condition in which an agent continues to operate with no traceable authorization path to a human principal. This section provides the formal vocabulary for that characterization, defines the constituent failure sub-modes, identifies the three structural conditions that enable them, catalogs four concrete failure modes, proves that depth limits alone are insufficient, and presents the formal model used in subsequent sections.

% -------------------------------------------------------------------------
\subsection{Formal Characterization}
\label{sec:drift-formal}

We model a multi-agent system as a rooted directed acyclic graph (DAG). Each node $n$ carries a \emph{parent pointer} $\alpha(n)$ — the node that authorized its provisioning. The root of every authorized chain is a human principal $h \in \mathcal{H}$. A \emph{spawn chain} $C = \langle n_0, n_1, \ldots, n_k \rangle$ satisfies $n_{i-1} = \alpha(n_i)$ for all $1 \leq i \leq k$ and $n_0 \in \mathcal{H}$.

\medskip
\noindent\textbf{Definition 1 (Orphaned Agent).}
An agent $n$ is \emph{orphaned} if its parent node $\alpha(n)$ satisfies one of:
\begin{enumerate}
  \item[(a)] \textbf{Terminated}: $\alpha(n)$ has exited operational state through completion, failure, or shutdown.
  \item[(b)] \textbf{Unreachable}: $\alpha(n)$ is non-responsive beyond the heartbeat timeout $\tau_{\mathrm{hb}}$.
\end{enumerate}
Orphaning severs the delegation chain at the parent. The agent persists in operational state but is structurally disconnected from its chain of authorization.

\medskip
\noindent\textbf{Definition 2 (Drifted Agent).}
An agent $n$ is \emph{drifted} if the chain $C(n) = \langle n_0, n_1, \ldots, n \rangle$ does not contain a human principal at its origin:
\begin{equation}
\mathrm{drifted}(n) \;\iff\; \neg\,\exists\, h \in \mathcal{H} : n_0 = h \text{ in the chain from } n \text{ to root}.
\label{eq:drifted-def}
\end{equation}
A drifted agent may have a living parent whose own chain does not reach any human principal. Drift is a property of the chain's origin, not of intermediate connectivity.

\medskip
\noindent\textbf{Definition 3 (Drift Triad).}
The \emph{drift triad} is the conjunction of three conditions that produces the RSP's target failure state:
\begin{equation}
\mathrm{DriftTriad}(n) = \langle \mathrm{orphaned}(n),\; \mathrm{drifted}(n),\; \mathrm{operating}(n) \rangle.
\label{eq:drift-triad}
\end{equation}
An agent enters the drift triad when it is simultaneously orphaned from its parent, unanchored to any human principal, and executing its event loop. The drift triad is the structural failure the Recursive Spawning Protocol is designed to make architecturally impossible.

\medskip
\noindent\textbf{Definition 4 (Drift Cascade).}
Drift is heritable: when a drifted agent $n$ spawns a child $c$, the condition propagates to all descendants:
\begin{equation}
\mathrm{drifted}(n, t) \land \mathrm{spawned}(n, c, t') \land t' > t \;\implies\; \mathrm{drifted}(c, t'')\quad \forall\, t'' > t'.
\label{eq:drift-cascade}
\end{equation}
A single undetected drift event at depth $d$ with branching factor $b$ produces up to $b^d$ drifted agents. This exponential amplification is what makes drift the central threat in recursive spawning systems.

% -------------------------------------------------------------------------
\subsection{Structural Conditions That Enable Drift}
\label{sec:drift-conditions}

Autonomous drift arises from three structural conditions endemic to conventional multi-agent runtimes. None is inherently malicious; each results from architectural choices that prioritize operational flexibility over accountability.

\medskip
\noindent\textbf{Condition 1 — Chain Opacity.}
In conventional runtimes, $\alpha(n)$ is a mutable reference. When a parent terminates or is reparented, the runtime updates $\alpha(n)$ to point to a new parent. This mutation retroactively rewrites the delegation chain: the child's current parent is accessible, but the original authorizer is not structurally retrievable. Without an immutable parent pointer established at provisioning, whether a chain reaches a human anchor cannot be answered at runtime.

\medskip
\noindent\textbf{Condition 2 — Scope Mutation.}
In conventional frameworks, an agent's operating scope $R(n)$ is re-evaluated at runtime by the agent itself, a supervisory node, or an external policy engine. When an agent encounters a condition outside its scope, the conventional response is scope expansion: the agent's authority is extended and recorded as a policy update, not as a Purpose Layer authorization requiring human principal approval. Scope mutation enables successive small expansions that collectively migrate the agent's operating scope far from its original authorization while the chain topology remains intact.

\medskip
\noindent\textbf{Condition 3 — Termination Ambiguity.}
In conventional frameworks, termination is event-driven: an agent terminates when it completes its task, receives a termination signal, or encounters an unrecoverable error. There is no structural requirement that termination be recorded in a tamper-evident ledger, that the agent's final state be audited, or that child nodes receive termination notification. A parent that exits cleanly without emitting termination events to its children leaves those children orphaned — they persist in operational state with no structural mechanism to detect parent disappearance.

% -------------------------------------------------------------------------
\subsection{Failure Modes}
\label{sec:drift-failure-modes}

The three enabling conditions produce four concrete failure modes, each a distinct pathway to the drift triad.

\medskip
\noindent\textbf{Failure Mode 1 — Silent Orphaning.}
A parent $\alpha(n)$ terminates without issuing a termination event to its child $n$. The child remains in \texttt{ACTIVE} state, continues to process tasks, may spawn further children, and has no knowledge its parent has exited. If the chain above $\alpha(n)$ does not reach a human anchor, $n$ is also drifted. Enabled by Condition 3.

\medskip
\noindent\textbf{Failure Mode 2 — Scope Creep Drift.}
A node $n$ encounters a condition $x \notin R(n)$ and does not escalate to its parent. Instead, $n$ expands its scope to include $x$ and continues operating. The expansion is recorded as a policy update, not a Purpose Layer authorization. The parent remains reachable and the chain topology is intact, but the effective operating scope of $n$ and all descendants has diverged from the human anchor's original intent. Successive expansions compound; each descendant inherits the expanded scope. Enabled by Condition 2.

\medskip
\noindent\textbf{Failure Mode 3 — Re-parenting Drift.}
A child $n$'s parent pointer $\alpha(n)$ is redirected to a new parent $p'$ by an administrative operation — for load balancing, organizational restructuring, or fault recovery. The redirection is logged as a configuration change, but the original authorization trace is not preserved. If $p'$ is not authorized by a human principal, $n$ becomes drifted through re-parenting without any change to its operational behavior. Enabled by Condition 1: the mutable parent pointer enables retroactive rewriting of the delegation chain — a form of \emph{accountability laundering}.

\medskip
\noindent\textbf{Failure Mode 4 — Ghost Authority.}
A node $n$ spawns a child $c$ without any authorized delegation chain — initiated by a bug, a compromised Bounds Engine, or an autonomous escalation that exceeded its authority. The child is provisioned without a valid parent pointer, without a human anchor reference, and without a Purpose Layer authorization record. The child enters operational state with a warrant that traces to no authorized chain. Enabled by the absence of a Fact Layer pre-commit hook.

% -------------------------------------------------------------------------
\subsection{Why Depth Limits Alone Do Not Solve Drift}
\label{sec:drift-depth-limits}

A common architectural response to unbounded agent spawning is a \emph{depth limit}: the system enforces a maximum chain depth $D_{\max}$, rejecting any spawn that would exceed this bound. Depth limits are necessary for preventing infinite spawn cascades but are \emph{insufficient} to prevent autonomous drift.

\medskip
\noindent\textbf{Theorem (Depth Limit Insufficiency).}
Let $D_{\max}$ be a system-enforced maximum chain depth. Then there exists at least one execution path in which a chain of depth $\leq D_{\max}$ produces a drifted agent, and no choice of $D_{\max}$ prevents this without additional structural constraints on scope refinement.

\medskip
\noindent\textbf{Proof.} Consider a spawn chain $C = \langle n_0, n_1, \ldots, n_k \rangle$ where $n_0$ is a human principal $h$ and $k < D_{\max}$. The depth constraint is satisfied by construction: $\mathrm{depth}(C) = k \leq D_{\max}$. Apply Failure Mode 2 (Scope Creep Drift): each node $n_i$ for $1 \leq i \leq k$ successively expands its operating scope $R(n_i)$ beyond $R(n_{i-1})$ without Purpose Layer authorization, recording each expansion as a policy update. After $k$ successive expansions, $R(n_k)$ is entirely disjoint from $R(n_0)$. The chain depth is $k \leq D_{\max}$; the chain has a human anchor at $n_0$; yet $n_k$ is drifted, because $R(n_k)$ is not a descendant refinement of $h(n_0)$'s intent. No choice of $D_{\max}$ prevents this, because the depth constraint operates on chain length, not on the scope relationship between successive nodes. $\blacksquare$

The drift prevention requirement is the conjunction of two independent constraints:
\begin{equation}
\mathrm{Drift\ Prevention} \;\iff\; \bigl( \mathrm{depth}(C) \leq D_{\max} \bigr)\ \land\ \bigl( \forall i:\ R(n_{i+1}) \subseteq R(n_i) \bigr).
\label{eq:drift-prevention-constraint}
\end{equation}
Depth limits enforce only the first conjunct. A system enforcing $D_{\max}$ without the scope refinement constraint is immune to unbounded spawn cascades but remains fully vulnerable to scope creep drift.

The deeper structural reason depth limits are insufficient is that they address the wrong dimension of the accountability problem. Depth limits constrain chain length; drift is fundamentally about authorization scope — whether each node's operating envelope is a legitimate descendant refinement of the human anchor's intent. A drifted agent at depth 1 with unbounded scope expansion authority can produce consequences exceeding the human anchor's intent as severely as a drifted agent at depth $D_{\max} - 1$.

% -------------------------------------------------------------------------
\subsection{Formal Model}
\label{sec:drift-model}

The system $\mathcal{S}$ is a 5-tuple:
\begin{equation}
\mathcal{S} = \langle \mathcal{N}, \alpha, R, D_{\max}, \mathcal{L} \rangle,
\label{eq:system-model}
\end{equation}
where $\mathcal{N}$ is the set of all agent nodes, $\alpha: \mathcal{N} \rightarrow \mathcal{N} \cup \{\bot\}$ is the parent pointer function ($\alpha(n) = \bot$ for the root human principal), $R: \mathcal{N} \rightarrow 2^{\mathcal{A}}$ maps each node to its operating scope, $D_{\max}$ is the maximum enforced chain depth, and $\mathcal{L}$ is the forensic ledger recording all authorized spawn events.

\medskip
\noindent\textbf{Chain depth} is defined recursively:
\begin{equation}
\mathrm{depth}(n) = 
\begin{cases}
0 & \text{if } \alpha(n) = \bot \quad (\text{root}) \\
1 + \mathrm{depth}(\alpha(n)) & \text{otherwise}.
\end{cases}
\label{eq:depth}
\end{equation}

\medskip
\noindent\textbf{Human anchor reachability} is:
\begin{equation}
\mathrm{chain-reaches}(n, h) \;\iff\; \exists\, \text{path } \langle n = n_k, n_{k-1}, \ldots, n_0 = h \rangle \text{ s.t. } \alpha(n_i) = n_{i-1}.
\label{eq:chain-reaches}
\end{equation}

\medskip
\noindent\textbf{Drift condition.} An agent $n$ is in the drift triad iff:
\begin{equation}
\bigl( \alpha(n) = \bot \ \lor \ \mathrm{unreachable}(\alpha(n)) \bigr)\ \land\ \bigl( \neg\,\exists\, h \in \mathcal{H} : \mathrm{chain-reaches}(n, h) \bigr)\ \land\ \mathrm{operating}(n) = \mathrm{true}.
\label{eq:drift-condition}
\end{equation}

\medskip
\noindent\textbf{Scope refinement invariant (RSP).} Under the Recursive Spawning Protocol, every spawn event from $n_i$ to $n_{i+1}$ must satisfy:
\begin{equation}
R(n_{i+1}) \subseteq R(n_i) \quad \text{and} \quad \mathrm{depth}(n_{i+1}) \leq D_{\max}.
\label{eq:scope-refinement}
\end{equation}
Both constraints are enforced by the Fact Layer pre-commit hook at provisioning time, making them structural invariants rather than policy conventions.

\medskip
\noindent\textbf{Drift propagation.} Without anchor verification at each spawn event:
\begin{equation}
\mathrm{drifted}(a) \;\implies\; \forall\, c \in \mathrm{children}(a):\ \mathrm{drifted}(c).
\label{eq:drift-propagation}
\end{equation}
Remediation requires explicit re-anchoring by a human actor: $\mathrm{re-establish\ chain-reaches}(a, h)$ for some $h \in \mathcal{H}$.

This formal model provides the vocabulary used throughout the remainder of this paper. The Recursive Spawning Protocol (Section~\ref{sec:rs-protocol}) modifies $\mathcal{S}$ by replacing the mutable parent pointer and mutable scope functions with their immutable counterparts, and by adding the Fact Layer pre-commit hook and forensic ledger that together enforce the constraints in Equation~\ref{eq:scope-refinement} as structural invariants — making the drift triad architecturally impossible rather than statistically unlikely.